\documentclass[12pt]{article}
\usepackage[margin=1in]{geometry}
\usepackage{setspace}
\usepackage{graphicx}
\usepackage{float}
\usepackage{hyperref}
\usepackage{enumitem}
\usepackage{booktabs}
\usepackage{longtable}

\onehalfspacing

\title{CMPE 321 Project 3: Modular DBMS Engine}
\author{Mert Ozustun, Selcen Celik}
\date{May 2026}

\begin{document}

\begin{titlepage}
    \centering
    \vspace*{3cm}

    {\LARGE \textbf{Project 3: Modular DBMS Engine}\par}
    \vspace{0.5cm}
    {\Large CmpE321 -- Introduction to Database Systems \par}
    {\Large Spring 2026\par}

    \vspace{2cm}

    {\large
    Mert Ozustun, 2022400192 \\
    Selcen Celik, 2022400219
    \par}

    \vspace{1cm}

    {\large Department of Computer Engineering \par}
    {\large Bogazici University \par}

    \vspace{1cm}
    {\large \today \par}
\end{titlepage}

\section{Introduction}

This report describes our modular DBMS engine. It has four layers.
From bottom to up they are the Disk Space Manager, the Buffer Manager, the
File and Index Manager, and the Query Processor. Each layer is a Python
package that interacts only with its lower layer. The entry point
\texttt{archive.py} reads a JSON config file, builds the four layers, and
feeds the input file line by line to the Query Processor.

The largest part of this report is Section~\ref{sec:sizing}. In that section, we describe every decision we made about sizing (integer width, string width, page header layout,
record format, catalog entry, index pages) and provide reasons for all of our desicions.
The instructor asked for this, so we made sure to keep it practical and provided 
calculations wherever needed. The subsequent sections explain each layer and then give the three
experiments with real numbers from our own runs.

\section{Architecture Overview}

The spec says the four layers are wired bottom-up in \texttt{archive.py}:

\begin{verbatim}
disk     = DiskSpaceManager(config)
buffer   = BufferManager(config, disk)
file_idx = FileIndexManager(config, buffer)
qp       = QueryProcessor(config, file_idx, buffer, disk)
\end{verbatim}

All calls between two layers produce a \emph{Result object}, and not a bare
value. The Result objects can be found in \texttt{models.py}:

\begin{itemize}[noitemsep]
  \item \texttt{PageResult} -- raw page bytes and if a real disk I/O
        occured. Returned by the Disk Space Manager.
  \item \texttt{BufferResult} -- a \texttt{PageResult} with cache-hit flag,
        evicted page id, and dirty-writeback flag. Returned by the Buffer
        Manager.
  \item \texttt{RecordResult} -- the matching records plus pages accessed,
        index nodes visited, records scanned, the strategy used, and a
        status. Returned by the File and Index Manager.
  \item \texttt{WriteResult} -- success flag, page id, and the bytes that
        were written. Returned by the Disk Space Manager on a write.
\end{itemize}

Since the data and the metadata always travel together within one of
these objects, a layer can change its internals and the layer above does
not have to change. For example, we can replace \texttt{LRU} with \texttt{MRU}
or \texttt{heap\_scan} with \texttt{bplus\_tree} by editing single line of
\texttt{config.json}. No code changes.

\section{Sizing and Formating Decisions}
\label{sec:sizing}

All numbers below are based on a default \texttt{page\_size} of 4096 bytes. The
header size is independent of the page size, so that changing the page size only
changes the body size.

\subsection{Field byte widths}

We record an integer field in \textbf{4 bytes}. We use Python
\texttt{struct} with the format \texttt{"<i"}. This corresponds to a signed 32-bit
little-endian integer. Range is about $-2.1\times10^{9}$ to
$+2.1\times10^{9}$. This is enough for normal data values and primary keys
in this project. We didn't use 8 bytes because that doubles the
integer cost for no benefit at the workload sizes the project asks for, and
it would lower how many records fit on a page.

We keep a string field in \textbf{32 bytes}, null-padded. If the value is
smaller than 32 bytes we pad the rest with \texttt{0x00}. If it is longer
we truncate it to 32 bytes. We choose 32 since the spec says string values are
alphanumeric only and the sample data (names like \texttt{Atreides},
\texttt{GiediPrime}) is short. 32 bytes fits those nicely and keeps the record small. Anyways, we want fixed width. Records
need to be fixed length so we can find slot $i$ by simple arithmetic instead of
scanning.

Both choices are fixed per build, so decoding a record is just a walk over
the field list adding 4 or 32 each step. No length prefixes are needed.

\subsection{Page header layout}

All pages (data, catalog, hash bucket, B+ tree node) have the same
32-byte header. The \texttt{struct} format is \texttt{"<BiI10si9s"}:

\begin{center}
\begin{tabular}{@{}rrl@{}}
\toprule
Offset & Size & Field \\
\midrule
0  & 1  & \texttt{page\_type} (0 data, 1 B+ internal, 2 B+ leaf, 3 hash bucket, 4 catalog) \\
1  & 4  & \texttt{page\_id} (signed) \\
5  & 4  & \texttt{record\_count} (unsigned; key count for index pages) \\
9  & 10 & \texttt{slot\_bitmap} (one byte per slot, 0 free / 1 used) \\
19 & 4  & \texttt{next\_page} (chained pages; $-1$ means none) \\
23 & 9  & reserved padding \\
\bottomrule
\end{tabular}
\end{center}

We have intentionally made the header a fixed size for all page types. It means
the body always starts at offset 32, so the slot math is the same
everywhere and the code is still short.

The slot bitmap is \textbf{one byte per slot}, not one bit per slot. A bitmap with at
most 10 slots is only 10 bytes either way, so bit-packing would
save 9 bytes per 4096-byte page (about 0.2\%) and make the slot
checks harder to read. The easier form of byte per slot was chosen. The 9
reserved bytes just pad the header out to a nice clean 32 so that future fields can
be added without shifting the body.

\subsection{Record format and page capacity}

The body of a data page is $4096 - 32 = 4064$ bytes. Records are packed
slot by slot beginning at offset 32. The size of a record is the sum of its
field widths.

Using the sample example, the \texttt{house} type with 3 string and 3 integer
fields:
\[
\text{record size} = 3\times 32 + 3\times 4 = 108 \text{ bytes}.
\]
By raw space a page could hold $4064 // 108 = 37$ such records, but the
config caps records per page at \texttt{max\_records\_per\_page} (at most
10). Thus we only store 10 and the rest of the body stays empty. In this case, the
limit is binding cap and not the byte space.

We examined the worst case as well. A type can have a maximum to 12 fields (see
below). The biggest record is 12 string fields:
\[
12 \times 32 = 384 \text{ bytes}, \qquad 10 \times 384 = 3840 \le 4064.
\]
Thus, the largest possible type still can fit 10 records in a single page. TThis explains why at 4096 bytes this is a 10 byte bitmap and setting a max cap of 10 is always safe. 
The cap can never exceed the volume the page can hold.

\subsection{System Catalog entry}

The catalog is stored in \texttt{\_\_catalog\_\_.bin} in the same
slotted-page format, so it survives restarts and every access goes through
the Buffer Manager. One type definition per slot in the catalog.

We don't keep a parsed copy of the catalog in memory. Each metadata
lookup requires a new scan of the catalog pages through the Buffer Manager. This keeps
us close to the project spec and means an unexpected crash cannot leave a
stale in-memory catalog behind; the only cache we rely on is the buffer
pool itself.

\begin{itemize}[noitemsep]
  \item Type name: \textbf{16 bytes}. The spec requires at least 12; we
        used 16 to give a little headroom and to keep the entry aligned.
  \item Field name: \textbf{24 bytes}. The spec reuires at least 20; we
        used 24 for the same reason.
  \item Fixed header part: \texttt{"<16siiiiiii"} $= 16 + 7\times4 = 44$
        bytes. The seven integers are number of fields, order of the primary key,
        number of data pages, number of records, B+ tree root page, number of hash buckets
        , and one reserved word.
  \item For each field descriptor: \texttt{"<24sB3s"} $= 24 + 1 + 3 = 28$
        bytes (name, a one-byte type code, three padding bytes).
\end{itemize}

A type can contaion between 6 and 12 fields. The minimum is 6 (spec
rule). We capped the upper bound at 12 so the catalog entry would be a fixed
size no matter how many fields the type really has:
\[
\text{catalog entry} = 44 + 12 \times 28 = 380 \text{ bytes}.
\]
With a body of 4064-byte that is  $4064 // 380 = 10$ entries per catalog page,
that fits with the 10-slot bitmap. When a catalog page is full we
link a new one via the header's \texttt{next\_page} field, so the
number of types is unbounded.

\subsection{Index page formats and fanout}

\textbf{Hash index} (\texttt{<type>\_idx.bin}). Pages $0..B-1$ are the
primary buckets and are allocated when the type is created. We default to 16
buckets. One bucket entry is \texttt{"<32sii"} $= 40$ bytes: a 32-byte key
(all keys serialized as a fixed 32-byte string so the bucket format is
the same for int and str keys), a 4-byte data page id, and a 4-byte slot
id. Each bucket page fits ten entries. When a bucket overflows we chain a new
page through \texttt{next\_page}. The hash function is the value mod the bucket
count for integers, and the sum of char codes mod the bucket count for
strings.

\textbf{B+ tree} (\texttt{<type>\_idx.bin}). Each node is one page. The key
is 4 bytes for an integer primary key and 32 bytes for a string primary
key, which is determined when the type is created. With a 4064-byte body the raw
fanout for an integer key is large (a leaf entry is key plus 8 bytes for
the data pointer, so about $4064 // 12 \approx 338$). To limit the fanout, we impose a hard limit
 \texttt{HARD\_FANOUT\_CAP} of 50 keys per node. Then a node never
needs more than one page, splits are cheap to write, and the tree
remains shallow. The three stays two or three levels deep with a fanout 
of 50 for the data sizes for the project, so a point lookup requires only  a few page reads.
Delete is lazy: we delete the entry from the leaf but do not rebalance. The
spec does not require rebalancing and a lazy delete is still correct, just
less compacted.

\subsection{Free space and page allocation}

Free space is provided on two levels. The slot bitmap within a page indicates
which slots are free. At the file level, each file keeps a small free list
of page ids that were deallocated and can be reused. We prefer a per-file
free list to a separate bitmap file because whole-page frees are rare in
this workload (a normal delete just clears a slot, it does not drop the
page), so a short list is sufficient and requires no extra file.

A page allocation is just bookkeeping. We increment the in-memory page count and
return the new id without writing anything to disk. The real bytes are
written later when the Buffer Manager flushes that page. Writing a zero
page at allocation time would mean one more disk write that the buffer's
later write-back makes pointless. A seek beyond the end of the file extends
it automatically if the first real write lands.

\subsection{I/O accounting decision}

On a write we do \emph{not} first read the old page.
\texttt{WriteResult.old\_data} exists for the Result-object contract but no
caller reads it, and pre-reading would incur one fake disk read for every
write. That will tie reads to writes and increase the read count. With our
choice a write is exactly one write. This makes the experiment numbers
honest and lets a write-heavy run actually show fewer reads than writes.

\section{The Four Layers}

\subsection{Layer 1 -- Disk Space Manager}

This is the only layer that touches files. It reads and writes fixed-size
pages with \texttt{seek()}; it never loads a whole file into memory. Each
relation, each index, and the catalog is a separate \texttt{.bin} file
beside \texttt{archive.py}. The path is derived from the script location, not
from the working directory and not from the config, which is what the spec
asks for.

It counts every read and every write and shows the totals. It also has a
\texttt{log\_write} stub that is called on every write. By default it does
nothing, but it is there and is called, as needed. At startup it scans the
folder for \texttt{.bin} files and recovers the number of pages in each file from the
file size. This is how state survives a restart.

\subsection{Layer 2 -- Buffer Manager}

The Buffer Manager is the page cache between Layer 3 and Layer 1. Layer 3
never calls the disk directly. The pool is an \texttt{OrderedDict} keyed by
(file id, page id). An \texttt{OrderedDict} provides O(1) move-to-end and
O(1) removal from either end, which is exactly what the two policies need.

On a hit there is no disk I/O and we move the page to the end to mark it as
recently used. On a miss we read the page from disk first (so a bad request
cannot evict a good page), then make room if the pool is full. To make
room we pick a victim by policy: \texttt{LRU} removes the page at the front
of the order (least recently used), \texttt{MRU} removes the page at the
back (most recently used). If the victim is dirty we write it back before
dropping it. Writes are deferred: \texttt{write\_page} just updates the
cached copy and marks it dirty. The actual disk write happens on eviction or
on the last \texttt{flush()}. We track requests, hits, misses, evictions,
and dirty writebacks.

\subsection{Layer 3 -- File and Index Manager}

This layer understands types, records, pages, and indexes. It reads the
catalog through the buffer, coerces input tokens into the declared field
types, and runs the operation with the current strategy.

\begin{itemize}[noitemsep]
  \item \texttt{heap\_scan}: scan every data page and check each used
        slot. This is the baseline and it is also stored under the two indexes
        (the indexes only store (page, slot) pointers into the same heap).
  \item \texttt{hash\_index}: a static hash on the primary key for equality
        lookups. The spec says heap scan is the fallback for a range query.
  \item \texttt{bplus\_tree}: a B+ tree on the primary key to support
        equality and range lookups. A range query on a non-primary-key
        field reverts to a heap scan.
\end{itemize}

The index is created at the time of the type creation and it is kept updated on
every insert and delete. Duplicate primary keys are rejected before
anything is changed. All index pages are handled by the Buffer Manager.

\subsection{Layer 4 -- Query Processor}

The Query Processor parses each line of input and passes it to Layer 3. It is responsible for 
the three output files. \texttt{output.txt} is cleared at the start of
a run and contains query results. The \texttt{stats\_output.txt} is overwritten on
every \texttt{stats} command since it is a snapshot, not a log.
The \texttt{log.csv} is append-only and is never cleared, so it survives
restarts. All operations are logged as \texttt{timestamp,operation,status}
where \texttt{int(time.time())} is used for the timestamp. Failures (duplicate key,
missing type, range on a non-integer field, garbage input) are logged as
\texttt{failure} and never cause the engine to crash.

Before creating the \texttt{stats} snapshot, the Query Processor calls
\texttt{buffer.flush()}. First, pages that are still dirty in the pool are written
, so they are counted as the disk writes they are about to become.
This guarantees that the reported totals are for the entire workload, not just the
I/O that happened to be on disk before the \texttt{stats} command was executed.

For \texttt{explain} we first query Layer 3 for an estimated strategy and
I/O. Then we pull the type's catalog entry into the buffer \emph{before}
snapshotting the counters, so the catalog access that is always present is not
counted as part of that one command. We snapshot the logical read and
write request counters, run the command, and report the difference as
``Actual I/O''. Hit and missdeltas of the buffer show how the cache served
those requests.

\section{Configuration, Operations, and Persistence}

In the config file, \texttt{page\_size}, \texttt{max\_records\_per\_page},
\texttt{buffer\_pool\_size}, \texttt{replacement\_policy}, and
\texttt{index\_strategy} are set up. The engine is run as
\texttt{python3 archive.py config.json input.txt} and only the config
changes between test runs.

Supported operations are create type, create record, delete record (by
primary key), search record (by primary key), range search (on any integer
field), \texttt{explain}, \texttt{stats}, and \texttt{stats reset}. All
data files, the catalog, the index files, and \texttt{log.csv} are
persistent. If the engine is stopped and then started again on the same folder,
it rebuilds its state from the \texttt{.bin} files and the catalog, and
queries return the old data without requiring \texttt{create type}. We have directly verified
this in our verification harness.

\section{Experiments}

All numbers below are from our own experiments. ``I/O'' is reads plus writes from
the \texttt{Disk I/O} line. ``Hit rate'' is the percentage on the
\texttt{Buffer Pool} line. The exact commands are in \texttt{record.txt} to
reproduced the runs. The workloads are generated by
\texttt{workload\_generator.py} with \texttt{--seed 42}.

\subsection{Experiment 1 -- LRU vs. MRU}

Setup: buffer pool size 8, index strategy \texttt{bplus\_tree}. The
sequential workload inserts 500 records and then does 50 full scans. The
random workload inserts 500 records and then does 200 equality searches.

\begin{table}[H]
\centering
\caption{LRU vs. MRU replacement policy comparison.}
\begin{tabular}{@{}lrrrr@{}}
\toprule
Workload & LRU I/Os & LRU Hit Rate & MRU I/Os & MRU Hit Rate \\
\midrule
Sequential & 37662 & 16.6\% & 35592 & 20.7\% \\
Random     & 15972 & 20.9\% & 13686 & 31.7\% \\
\bottomrule
\end{tabular}
\end{table}

MRU was better than LRU in both lines. The sequential row is the
sequential flooding problem. The full scan hits more pages than the pool
can contain. With LRU, the scan finishes and the pool is filled with the
latest pages.The next scan starts again at the first page and that is the
page LRU just threw away. So scan after scan, almost every page is a miss. MRU discards the most recently used page and retains the older ones, 
so when the next scan begins some early pages are still in the pool and are hits. 
Therefore, MRU has the higher hit rate on
repeated scans.

The random row is the same for a related reason. Every B+ tree
search begins at the root and the upper internal nodes. Those get loaded
early and, remain in the pool on MRU, while the leaf pages
churning at the most-recent end and get evicted. Under LRU the root may drift
towards the eviction end between searches and be dropped, so the next
search has to read it again. MRU’s higher hit rate here is due to retaining the
few hot upper nodes.

\subsection{Experiment 2 -- Index Strategy Comparison}

Setup: buffer pool size 16, replacement policy \texttt{LRU}. Equality works
300 records and 150 random equality searches. Range is based on 300 records and
150 random range queries on an integer field.

\begin{table}[H]
\centering
\caption{I/O comparison across index strategies (reads + writes).}
\begin{tabular}{@{}lrrr@{}}
\toprule
Query Type & HeapScan & HashIndex & B+-Tree \\
\midrule
Equality & 9121  & 5393                & 5148 \\
Range    & 11795 & 9759 (heap fallback) & 11019 \\
\bottomrule
\end{tabular}
\end{table}

As expected, the heap scan is the worst for equality, reading every data page until
it finds the key.The hash index and the B+ tree go almost straight to the record, so do much less I/O. The B+ tree is slightly
better than the hash index here since the hash buckets have some overflow
chaining while the tree stays shallow.

The heap scan reads everything again for the range query. The spec requires 
the hash index to fall back to a heap scan for ranges, so it is doing the 
same kind of full scan, and its lower number compared to plain heap is from a
different buffer state left over from the insert phase (building the hash index 
touches different pages than building nothing), not from a better scan. The 
B+ tree does not save much here because the ranges generated are wide, 
often a large part of the key space. Still a lot visits most leaves and 
then follows pointers into most data pages plus the index pages on top.
The B+ tree only really wins on selective ranges. It is still at or below 
the heap scan, which matches what we expect.

\subsection{Experiment 3 -- Buffer Pool Size Sensitivity}

Setup: random workload (400 records, 200 equality searches), index
strategy \texttt{bplus\_tree}, replacement policy \texttt{LRU}. Only the
buffer pool size is changed.

\begin{table}[H]
\centering
\caption{Effect of buffer pool size on I/O performance.}
\begin{tabular}{@{}lrr@{}}
\toprule
Buffer Size & I/Os & Hit Rate \\
\midrule
4  & 11023 & 23.5\% \\
8  & 10717 & 25.2\% \\
16 & 9481  & 33.1\% \\
32 & 5079  & 63.8\% \\
64 & 54    & 100.0\% \\
\bottomrule
\end{tabular}
\end{table}

The pattern is clean: more frames means less missed reads and less I/O. As the
pool grows it can hold more of the working set (the upper B+ tree nodes and
the hot data pages), so searches hit instead of miss. The interesting point is the
jump at 64. The relation haas only 40 data pages,plus a small number of B+ tree pages.
There are 400 records at 108 bytes and 10 records per
page. The entire dataset fits in a pool of 64. After the initial warmup
all requests are hits and the only disk I/O remaining is the 54 writes from
the final flush of the dirty pages. Reads fall to zero and the hit rate is
100\%. This also showing the write-back design working: nothing is written to
disk until it has to be.

\section{Conclusion}

The engine matches the spec. The four layers are independent and only talk
through Result objects, so the policy and the index strategy can be swapped
from the config with no code change. We recorded each size we selected and
showed the math that proves a page always contains its 10 records even
for the largest allowed type. The three experiments behave as the
theory predicts: MRU beats LRU under sequential flooding, the indexes beat
the heap scan for equality, and a larger buffer pool steadily lowers I/O
until the data fits and disk reads disappear.

\end{document}
